%*******************************************************************************
% Macro per il documento corrente
%*******************************************************************************
\newcommand{\sharedPath}{../shared}
\newcommand{\doctitle}{Business Plan}

%*******************************************************************************
% Preambolo
%*******************************************************************************
\documentclass[a4paper,10pt,twoside]{book}

%*******************************************************************************
% Codifica e lingua
%*******************************************************************************
\usepackage[utf8x]{inputenc}
\usepackage[T1]{fontenc}
\usepackage[english,italian]{babel}
\usepackage{lmodern}
\usepackage{multirow}
\usepackage{amsmath,eurosym}
%*******************************************************************************
% Qualche macro utile a tutti
%*******************************************************************************
\newcommand{\docRoot}{..}
\newcommand{\inglese}[1]{\foreignlanguage{english}{\textit{#1}}}
\newcommand{\team}{\textsf{EtaBeta\,Software}\xspace}
\newcommand{\caName}{BPM-1.0\xspace}
\newcommand{\customer}{\textsf{Alpha\,\&\,Partners}\xspace}
\newcommand{\mktg}{\inglese{marketing}\xspace}
\newcommand{\bsn}{\inglese{business}\xspace}
\newcommand{\sw}{\inglese{software}\xspace}

%*******************************************************************************
% Figure e immagini
%*******************************************************************************
\usepackage{graphicx}
\graphicspath{{\docRoot/shared/pictures/}}

%*******************************************************************************
% Tabelle
%*******************************************************************************
\usepackage{booktabs}
\usepackage{array}

%*******************************************************************************
% Elenchi puntati personalizzati
%*******************************************************************************
\usepackage{enumitem}

%*******************************************************************************
% Landscape per ruotare le pagine
%*******************************************************************************
\usepackage{pdflscape}

%*************************************************
% Collegamenti intra- e intertestuali
%*************************************************
\usepackage{hyperref}
\hypersetup{%
    colorlinks=false,linktocpage=false,pdfborder={0 0 0},%
    pdfstartpage=1, pdfstartview=FitV,plainpages=false,%
    urlcolor=Black, linkcolor=Black,
    pdfcreator={pdfLaTeX},%
    pdfproducer={pdfLaTeX with hyperref package}%
}

%*************************************************
% Grafica vettoriale
%*************************************************
\usepackage{tikz}
\usetikzlibrary{shadows,arrows,shapes}

%*******************************************************************************
% Altri pacchetti
%*******************************************************************************
\usepackage{xspace} % per spazi condizionali extra
\usepackage{lastpage} % per sapere il numero totale di pagine
\usepackage{microtype} % qualche accorgimento tipografico

% *************************************************
% Intestazioni e piè di pagina personalizzati
% *************************************************
\usepackage{fancyhdr}

% stile normale
\fancypagestyle{normal}{
\fancyhead{}
\fancyhead[RE,LO]{
\sffamily\team
}
\fancyhead[LE,RO]{
\sffamily\leftmark
}
\renewcommand{\headrulewidth}{.4pt}
\cfoot{}
\fancyfoot[LE,RO]{\sffamily
  pag. \thepage{} di \pageref{LastPage}}
\fancyfoot[RE,LO]{\sffamily\doctitle}
\renewcommand{\footrulewidth}{.4pt}
}

% stile per gli indici
\fancypagestyle{toc}{
\fancyhead{}
\renewcommand{\headrulewidth}{0pt}
\cfoot{}
\fancyfoot[RO,LE]{\sffamily\thepage{}}
\fancyfoot[RE,LO]{\sffamily\doctitle{}}
\renewcommand{\footrulewidth}{.4pt}
}

\pagestyle{fancy}
\renewcommand{\sectionmark}[1]{\markboth{#1}{#1}}


%*******************************************************************************
% Inizio documento
%*******************************************************************************
\begin{document}

\pagestyle{empty}
\begin{center}

{\sffamily
Sviluppo e Gestione Progetti\\
a.a. 2012--2013
}

\vskip 1.5cm

{
\setlength{\fboxsep}{.2pt}
\fbox{\includegraphics[width=\textwidth]{logo}}
}

\vskip .5cm

{\Huge\sffamily\bfseries
\team
}

\vskip 1.5cm

% titolo del progetto
{\Large\sffamily\bfseries
\caName
}

\vskip 1cm

% titolo del documento
\hrule
\vskip 10pt
{\Huge\scshape
\doctitle
}
\vskip 10pt
\hrule

\end{center}

\clearpage
\pagenumbering{roman} 

\tableofcontents{\thispagestyle{toc}}

\clearpage

\pagestyle{normal}
\pagenumbering{arabic}


\section{Introduzione}
\subsection{Scopo del documento}
Il presente documento si propone di fornire informazioni su \team, descrivendo quali sono le abilità e le competenze che intende mettere in campo per la realizzazione del progetto commissionato da \customer delineando una \inglese{roadmap} di sviluppo attraverso cui raggiungere gli obiettivi di progetto (sezione~\ref{sec:whoweare}).

In seguito, verranno presentati con maggiori dettagli gli obiettivi del progetto, i problemi che la soluzione proposta si propone di affrontare e le funzionalità che quest'ultima dovrà soddisfare tramite l'illustrazione degli scenari tipici di utilizzo (sezione \ref{sec:whattheproblemis}).
 
Infine, sarà illustrata la strategia di uscita, descrivendo con maggiori dettagli perché la soluzione proposta risolve il problema, le strategie di \mktg e quali possono essere gli utilizzi futuri e i potenziali acquirenti esterni (sezione \ref{sec:exitstrategy}).

\subsection{Scopo del progetto}
%TODO da terminare!

\subsection{Executive summary}
%TODO questo va scritto alla fine, non si può fuffare ora!

\section{L'impresa}\label{sec:whoweare}
% idea di impresa -> forse simile allo scopo del progetto
% risorse: tecnologia, brevetti, prodotti&servizi
% competenze (background di formazione del team di sviluppo)
% successi della storia aziendale (traction)
% modello business e quanto verrà a costare (o vengono a costare progetti simili)
% obiettivi raggiunti e da raggiungere nel corso del tempo (roadmap)

\subsection{Analisi competitiva}
%TODO fare analisi SWOT e altre cazzate!

\section{Il mercato di riferimento}\label{sec:whattheproblemis}
\subsection{Analisi di mercato}
%  - analisi della domanda
%  - analisi del settore
\subsection{Analisi del problema}
%  - problem breakdown
%  - obiettivi e milestones

\section{Strategie di \mktg}\label{sec:exitstrategy}
% tecniche di mktg (eventualmente riferendo progetti già esistenti)
% piano di mktg
% piano commerciale

\section{Analisi finanziaria}
% prospettive di vendite e guadagni rispetto alle spese sostenute (almeno su 3 anni)
% capitale come grafico in prospettiva

\end{document}